% ============================================================================
%  Capitolo 3 -- Storia della computazione                       (Lezione 03)
% ============================================================================
\chapter{Storia della computazione}
\label{ch:history}

\epigraphline{\texttt{printf("Hello, world!");}}
  {Brian Kernighan e Dennis Ritchie, \emph{The C Programming Language}, 1978}

\begin{chapgoals}
Alla fine di questo capitolo saprete
\begin{itemize}
\item definire l'informatica e dire perché non è la scienza dei calcolatori;
\item collocare su una linea del tempo le macchine e le persone chiave,
  dall'abaco al calcolatore a programma memorizzato;
\item nominare le quattro unità dell'architettura di von Neumann e ritrovarle
  nella macchina programmata nel Capitolo~\ref{ch:firstcontact};
\item spiegare che cosa significa RISC e perché \riscv è un insieme di
  istruzioni aperto;
\item descrivere dove si colloca l'informatica tra scienza, ingegneria ed economia.
\end{itemize}
\end{chapgoals}

Il lettore ha già eseguito un programma \riscv. Questo capitolo racconta da
dove viene quella macchina. Non è una cronaca: delle migliaia di macchine e di
persone che la storia della computazione potrebbe elencare, conserva quelle che
hanno aggiunto un'idea che si trova ancora dentro l'ESP32-C6 sulla breadboard.
Il filo da seguire è che, a ogni tappa, qualcuno ha trovato un modo nuovo di
fare una di due cose: \emph{rappresentare} l'informazione in un oggetto fisico,
oppure \emph{prescrivere} una sequenza di operazioni da eseguire senza doverci
più pensare.

% ----------------------------------------------------------------------------
\section{Che cos'è l'informatica?}

\begin{definition}{Informatica}
L'informatica è la scienza applicata che si occupa del trattamento
dell'informazione mediante procedure automatizzate.
\end{definition}

La parola \emph{informatica} viene dal francese \emph{informatique},
contrazione di \emph{information automatique}. L'inglese dice \emph{computer
science} e, sempre più spesso, \emph{informatics}. Lo strumento della
disciplina si chiama \emph{elaboratore} o \emph{calcolatore} in italiano,
\emph{ordinateur} in francese, \emph{computer} in inglese. Il nome inglese
della disciplina ha il difetto di mettere lo strumento al posto dell'oggetto di
studio, un difetto che una frase attribuita a Edsger W.\ Dijkstra corregge:

\begin{quote}\itshape
L'informatica non riguarda i calcolatori più di quanto l'astronomia riguardi i
telescopi.
\end{quote}

La macchina è lo strumento, non l'oggetto di studio. In questo libro lo
strumento viene studiato da vicino proprio perché smetta di essere un mistero.

La disciplina ha due anime. L'\emph{informatica teorica} studia i fondamenti
dell'informazione e della computazione a livello logico: che cosa si può
calcolare, e a quale costo. Il suo oggetto centrale è l'\term{algoritmo}, un
insieme finito di operazioni che descrive la trasformazione di un dato $A$ in
un risultato $B$. L'\emph{informatica applicata} studia le tecnologie
abilitanti e i sistemi informativi costruiti con esse. È pervasiva, è un motore
dello sviluppo economico e dell'intero settore delle tecnologie
dell'informazione e della comunicazione, è una disciplina STEM al centro di una
rivoluzione industriale; e ha un costo sociale che va governato, il
\term{divario digitale} tra chi ha accesso a queste tecnologie e chi non lo ha.

% ----------------------------------------------------------------------------
\section{Prima dell'elettricità}

Da tremila anni si costruiscono macchine che aiutano a calcolare.

\subsection{Le origini: l'abaco e il meccanismo di Anticitera}

Babilonesi, Cinesi, Greci e Romani usavano tutti dei pallottolieri
(Figura~\ref{fig:hist:origins}). In un abaco un numero \emph{è} una posizione
di palline, e un'addizione \emph{è} un movimento delle dita: il dispositivo è
una rappresentazione fisica di dati e di operazioni. Il meccanismo di
Anticitera, recuperato da un relitto e costruito tra il secondo e il primo
secolo avanti Cristo (la maggior parte delle stime cade tra il 205 e il
60~a.C. circa), è il più antico calcolatore a ingranaggi conosciuto: i suoi
quadranti calcolavano posizioni astronomiche ed eclissi, e in esso un numero è
l'angolo di una ruota dentata.

\begin{figure}[htb]
\centering
\begin{subfigure}[b]{0.24\textwidth}\includegraphics[width=\textwidth,height=32mm,keepaspectratio]{abacus_roman.jpg}\caption{Abaco romano tascabile}\end{subfigure}\hfill
\begin{subfigure}[b]{0.24\textwidth}\includegraphics[width=\textwidth,height=32mm,keepaspectratio]{abacus_chinese.jpg}\caption{\emph{Suanpan} cinese}\end{subfigure}\hfill
\begin{subfigure}[b]{0.24\textwidth}\includegraphics[width=\textwidth,height=32mm,keepaspectratio]{antikythera.jpg}\caption{Meccanismo di Anticitera}\end{subfigure}\hfill
\begin{subfigure}[b]{0.24\textwidth}\includegraphics[width=\textwidth,height=32mm,keepaspectratio]{antikythera_plate.jpg}\caption{Una delle sue lastre}\end{subfigure}
\caption{Le origini: un numero diventa una posizione di palline o l'angolo di un ingranaggio.}
\label{fig:hist:origins}
\end{figure}

\begin{prgconcept}
La rappresentazione viene prima del calcolo. Prima di poter calcolare
qualcosa, un oggetto fisico deve stare per un numero.
\end{prgconcept}

\subsection{Il mondo islamico medievale}

Tra il nono e il tredicesimo secolo gli studiosi del mondo islamico hanno
conservato la scienza dell'antichità e vi hanno aggiunto tre cose che
riguardano questa storia. La prima è una famiglia di strumenti astronomici,
l'astrolabio meccanico (Figura~\ref{fig:hist:astrolabe}) e il torquetum, che
sono calcolatori analogici per le posizioni degli astri. La seconda è la
crittoanalisi: nel nono secolo al-Kindi descrisse l'analisi delle frequenze, il
metodo per violare un cifrario contando quante volte compare ciascun simbolo.
La terza è la macchina programmabile: il suonatore di flauto automatico dei
fratelli Ban\=u M\=us\=a, e gli orologi e gli automi a pesi di al-Jazari
(1206), il cui comportamento si poteva cambiare spostando dei pioli.

La stessa parola \emph{algoritmo} viene dal nome del matematico
al-Khw\=arizm\=\i, come \emph{algebra} viene dal titolo del suo libro.

\begin{figure}[htb]
\centering
\begin{subfigure}[b]{0.3\textwidth}\centering\includegraphics[width=\textwidth,height=42mm,keepaspectratio]{astrolabe.jpg}\caption{Un astrolabio di ottone}\label{fig:hist:astrolabe}\end{subfigure}\hfill
\begin{subfigure}[b]{0.3\textwidth}\centering\includegraphics[width=\textwidth,height=42mm,keepaspectratio]{pascaline.jpg}\caption{Una Pascalina}\label{fig:hist:pascaline}\end{subfigure}\hfill
\begin{subfigure}[b]{0.3\textwidth}\centering\includegraphics[width=\textwidth,height=42mm,keepaspectratio]{leibniz_reckoner.jpg}\caption{La calcolatrice di Leibniz}\end{subfigure}
\caption{Strumenti e calcolatrici prima del 1700.}
\end{figure}

\subsection{Dal 1600 al 1800: calcolatrici meccaniche e logica}

Il Seicento costruì le prime calcolatrici meccaniche. Le tavole dei logaritmi
di Nepero (1614) trasformarono le moltiplicazioni in addizioni. La
\emph{Pascalina} di Blaise Pascal (1642) sommava e sottraeva, e la sua vera
invenzione fu il \term{riporto} automatico: ogni ruota è una cifra decimale
(Figura~\ref{fig:hist:pascaline}), e una ruota che passa da 9 a 0 spinge avanti
di uno scatto la vicina, esattamente come farà con i bit il sommatore completo
del Capitolo~\ref{ch:combinational}. La calcolatrice a cilindri scalati di
Leibniz (1673) eseguiva le quattro operazioni.

Lo stesso Leibniz descrisse l'aritmetica binaria nel 1703, e un secolo e mezzo
più tardi George Boole, in \emph{The Laws of Thought} (1854), mostrò che la
logica si può trattare come un'algebra \citep{boole1854}: l'argomento dei
Capitoli \ref{ch:boole1} e~\ref{ch:boole2}.

Nel frattempo una macchina che non aveva nulla a che fare con i numeri
introdusse il programma memorizzato. Nel telaio di Jacquard (1804) il disegno
del tessuto è registrato su una catena di schede perforate: si cambiano le
schede, e lo stesso telaio tesse un disegno diverso.

\subsection{L'Ottocento: Charles Babbage e Ada Lovelace}

Charles Babbage unì le due linee (Figura~\ref{fig:hist:babbage}). La sua
\emph{Difference Engine}, la macchina alle differenze, calcolava
automaticamente tavole matematiche. La sua \emph{Analytical Engine}, la
macchina analitica, progettata ma mai completata, era un calcolatore di uso
generale in tutto tranne che nella tecnologia: aveva un \emph{mill}, il mulino,
cioè l'unità aritmetica; un'unità di controllo; uno \emph{store}, il magazzino,
cioè la memoria; e un dispositivo di uscita. Programmi e dati dovevano essere
forniti su schede perforate, prese in prestito dal telaio di Jacquard.

Ada Lovelace ideò un metodo per programmare la macchina analitica. Le sue note
del 1843 contengono quello che è considerato il primo programma pubblicato, per
il calcolo dei numeri di Bernoulli, e lei è considerata la prima programmatrice
della storia.

\begin{figure}[htb]
\centering
\begin{subfigure}[b]{0.32\textwidth}\centering\includegraphics[width=\textwidth,height=40mm,keepaspectratio]{difference_engine.jpg}\caption{Difference Engine}\end{subfigure}\hfill
\begin{subfigure}[b]{0.32\textwidth}\centering\includegraphics[width=\textwidth,height=40mm,keepaspectratio]{analytical_engine.jpg}\caption{Analytical Engine}\end{subfigure}\hfill
\begin{subfigure}[b]{0.32\textwidth}\centering\includegraphics[width=\textwidth,height=40mm,keepaspectratio]{punched_cards.jpg}\caption{Schede perforate}\end{subfigure}
\caption{Le macchine di Charles Babbage.}
\label{fig:hist:babbage}
\end{figure}

% ----------------------------------------------------------------------------
\section{Il Novecento}

Il Novecento ha prodotto l'idea di una macchina universale, e poi le macchine
che l'hanno resa reale. La teoria è venuta prima.

\subsection{Turing, Shannon e Nakajima}

Nel 1936 Alan Turing definì una macchina astratta universale, e con essa la
nozione di ciò che si può calcolare \citep{turing1936}. Durante la guerra
lavorò a Bletchley Park alla violazione del cifrario tedesco Enigma
(Figura~\ref{fig:hist:enigma}).

Negli stessi anni Akira Nakajima in Giappone e Claude Shannon negli Stati
Uniti fondarono la teoria dei circuiti a commutazione: l'algebra di Boole
descrive i circuiti fatti di interruttori \citep{shannon1938}. È il principio
su cui poggia tutta la progettazione digitale, e il ponte tra la logica dei
Capitoli \ref{ch:boole1}--\ref{ch:boole2} e i circuiti dei Capitoli
\ref{ch:electronics}--\ref{ch:combinational}.

In Germania Konrad Zuse costruì lo Z3 (1941), il primo calcolatore funzionante
controllato da un programma, di cui nel 1998 è stata dimostrata la completezza
secondo Turing; il suo successore, lo Z4, fu uno dei primi calcolatori venduti
commercialmente.

\begin{figure}[htb]
\centering
\begin{subfigure}[b]{0.32\textwidth}\centering\includegraphics[width=\textwidth,height=38mm,keepaspectratio]{enigma.jpg}\caption{Una macchina Enigma}\label{fig:hist:enigma}\end{subfigure}\hfill
\begin{subfigure}[b]{0.32\textwidth}\centering\includegraphics[width=\textwidth,height=38mm,keepaspectratio]{zuse_z3.jpg}\caption{Zuse Z3 (replica)}\end{subfigure}\hfill
\begin{subfigure}[b]{0.32\textwidth}\centering\includegraphics[width=\textwidth,height=38mm,keepaspectratio]{colossus.jpg}\caption{Colossus}\end{subfigure}
\caption{Macchine degli anni Quaranta.}
\end{figure}

\subsection{John von Neumann: il calcolatore a programma memorizzato}

Nel 1945 il rapporto sull'EDVAC, firmato da John von Neumann, descrisse
un'architettura hardware per calcolatori digitali programmabili
\citep{vonneumann1945}. Il suo cuore è un'idea sola.

\begin{prgconcept}[Il programma memorizzato]
Programma e dati condividono la stessa memoria: \emph{un programma è un dato}.
Le istruzioni sono numeri conservati in memoria, prelevati ed eseguiti uno dopo
l'altro.
\end{prgconcept}

L'architettura ha quattro unità essenziali (Figura~\ref{fig:hist:vn}):
un'\term{unità di controllo}, che governa la sequenza delle operazioni;
un'\term{unità aritmetico-logica} (ALU), con il suo registro accumulatore; una
\term{memoria}; e dispositivi di \term{ingresso} e di \term{uscita}.

\begin{figure}[htb]
\centering
\begin{tikzpicture}[node distance=6mm and 8mm]
  \node[blk,fill=prgsteel,text=white,minimum width=62mm,minimum height=11mm] (m) {\textbf{Memoria}: programma e dati};
  \node[blk,fill=prgaccent,text=white,minimum width=27mm,minimum height=15mm,below left=10mm and -27mm of m] (cu) {\textbf{Unità di controllo}};
  \node[blk,fill=prgamber,text=white,minimum width=27mm,minimum height=15mm,below right=10mm and -27mm of m] (alu) {\textbf{ALU}\\(accumulatore)};
  \node[blk,fill=prggreen,text=white,minimum width=27mm,below=10mm of cu] (in) {\textbf{Ingresso}};
  \node[blk,fill=prggreen,text=white,minimum width=27mm,below=10mm of alu] (out) {\textbf{Uscita}};
  \draw[<->,thick] (cu.north) -- (cu.north|-m.south);
  \draw[<->,thick] (alu.north) -- (alu.north|-m.south);
  \draw[<->,thick] (cu) -- (alu);
  \draw[->,thick] (in) -- (cu);
  \draw[->,thick] (alu) -- (out);
\end{tikzpicture}
\caption{Le quattro unità dell'architettura di von Neumann.}
\label{fig:hist:vn}
\end{figure}

\begin{prgmachine}
L'ESP32-C6 è ancora una macchina di von Neumann, e le sue quattro unità erano
tutte al lavoro nel Capitolo~\ref{ch:firstcontact}. I registri e la ALU hanno
eseguito \code{addi t1, t1, 1}; l'unità di controllo ha deciso che l'istruzione
successiva era \ins{sw}; la memoria conteneva la variabile \code{counter}
\emph{e le istruzioni stesse}, tra cui il numero \hex{00130313}; il joystick e
il display sono ingresso e uscita.
\end{prgmachine}

\subsection{Otto decenni di calcolatori}

Da allora ogni decennio ha reso il calcolatore più piccolo, più economico e più
vicino alle persone (Tabella~\ref{tab:hist:decades},
Figura~\ref{fig:hist:decades}); nessuno ha cambiato l'idea del programma
memorizzato.

\begin{table}[htb]
\centering\small
\caption{Otto decenni di calcolatori.}\label{tab:hist:decades}
\begin{tabular}{@{}lll@{}}
\toprule
\textbf{Decennio} & \textbf{Epoca} & \textbf{Che cosa è cambiato}\\
\midrule
Anni '40 & La guerra & ENIAC (USA) e Colossus (Regno Unito), 1943--45\\
Anni '50 & Militari e governi & EDVAC, 1949: il programma memorizzato\\
Anni '60 & I mainframe & una macchina, molti utenti\\
Anni '70 & I minicomputer & un calcolatore per reparto\\
Anni '80 & Il personal computer & un calcolatore per persona\\
Anni '90 & Internet & i calcolatori si collegano\\
Anni 2000 & Smartphone e tablet & un calcolatore in tasca\\
Anni 2010 & Il cloud computing & il calcolo come servizio\\
Anni 2020 & Internet delle cose e IA & calcolatori ovunque\\
\bottomrule
\end{tabular}
\end{table}

\begin{figure}[htb]
\centering
\begin{subfigure}[b]{0.24\textwidth}\centering\includegraphics[width=\textwidth,height=28mm,keepaspectratio]{eniac.png}\caption{ENIAC}\end{subfigure}\hfill
\begin{subfigure}[b]{0.24\textwidth}\centering\includegraphics[width=\textwidth,height=28mm,keepaspectratio]{mainframe.jpg}\caption{Un mainframe}\end{subfigure}\hfill
\begin{subfigure}[b]{0.24\textwidth}\centering\includegraphics[width=\textwidth,height=28mm,keepaspectratio]{minicomputer.jpg}\caption{Un minicomputer}\end{subfigure}\hfill
\begin{subfigure}[b]{0.24\textwidth}\centering\includegraphics[width=\textwidth,height=28mm,keepaspectratio]{c64.jpg}\caption{Commodore 64}\end{subfigure}
\caption{Da una stanza a una scrivania in quarant'anni.}
\label{fig:hist:decades}
\end{figure}

% ----------------------------------------------------------------------------
\section{Dagli home computer a RISC-V}

\subsection{I microprocessori: da CISC a RISC a uno standard aperto}

Il processore del kit \prg è fatto com'è per una sequenza di decisioni prese
nell'arco di cinquant'anni (Tabella~\ref{tab:hist:cpu}). I primi
microprocessori sono cresciuti per accumulo: ogni generazione aggiungeva
istruzioni, spesso complicate, per facilitare il lavoro dei compilatori e dei
programmatori in assembly. Il risultato si chiama \term{complex instruction set
computer} (CISC), e la famiglia x86, nata nel 1978, porta ancora questa eredità
nei personal computer di oggi.

All'inizio degli anni Ottanta due progetti universitari, a Berkeley e a
Stanford, misurarono che cosa fanno davvero i programmi, e scoprirono che un
processore con istruzioni poche, semplici e regolari può essere reso più veloce
di uno con molte istruzioni complicate. È l'idea del \term{reduced instruction
set computer} (RISC), per la quale David Patterson e John Hennessy hanno
ricevuto il premio Turing nel 2017. ARM l'ha applicata al basso consumo, e i
core ARM sono oggi in quasi ogni telefono. Nel 2010 il quinto progetto RISC di
Berkeley è diventato qualcosa di nuovo: un insieme di istruzioni pubblicato
come \emph{standard aperto}.

\begin{table}[htb]
\centering\small
\caption{I microprocessori: da CISC a RISC a uno standard aperto.}\label{tab:hist:cpu}
\begin{tabularx}{\textwidth}{@{}lP{0.31\textwidth}L@{}}
\toprule
\textbf{Anno} & \textbf{Tappa} & \textbf{Perché conta}\\
\midrule
1971 & Intel 4004 & il primo microprocessore commerciale: 2\,300 transistor su un chip\\
1975 & MOS 6502 & una CPU a 8 bit economica: Apple II, Commodore 64, l'era degli home computer\\
1978 & Intel 8086 & l'inizio di x86, un insieme di istruzioni complesso (CISC) ancora nei PC di oggi\\
1980--82 & RISC di Berkeley e MIPS di Stanford & istruzioni poche, semplici e regolari vanno più veloci: l'idea RISC\\
1985 & ARM1 & RISC per il basso consumo: oggi in quasi ogni telefono\\
2010 & \riscv nasce a UC Berkeley & il quinto progetto RISC di Berkeley diventa un insieme di istruzioni aperto\\
2015--20 & RISC-V Foundation, poi RISC-V International & una specifica condivisa, governata dalla Svizzera\\
2023--25 & European Chips Act; il progetto DARE & \riscv come strategia per la sovranità tecnologica europea\\
\bottomrule
\end{tabularx}
\end{table}

\begin{definition}{RISC-V}
\riscv è una \term{instruction set architecture}, un'architettura dell'insieme
di istruzioni, a standard aperto: la specifica delle istruzioni, dei registri e
della loro codifica è pubblica, e chiunque può progettare, costruire e vendere
un processore che la realizza senza pagare una licenza per l'architettura.
\end{definition}

\riscv quindi non è un processore, e non è stato il primo RISC: il RISC di
Berkeley, MIPS e ARM sono venuti prima. Ciò che lo distingue da x86 e da ARM è
l'apertura.

\subsection{La legge di Moore in un numero}

L'Intel 4004 del 1971 aveva 2\,300 transistor. Un processore per portatili di
oggi ne ha decine di miliardi: una crescita di un fattore di circa dieci
milioni. Per cinquant'anni il numero di transistor su un chip è raddoppiato
circa ogni due anni, un'osservazione fatta da Gordon Moore nel 1965 e rivista
nel 1975 \citep{moore1965}. È un'osservazione su un'industria, non una legge di
natura, e i suoi limiti fisici sono ormai vicini. Il progresso viene sempre più
da architetture migliori, da acceleratori specializzati e da progetti aperti,
ed è per questo che l'architettura, l'argomento di questo libro, conta più che
mai. Il Capitolo~\ref{ch:beyond} riprende la storia da qui.

\subsection{Lo spirito degli home computer, su una CPU moderna e aperta}

Gli home computer degli anni Ottanta avevano una qualità che le macchine
successive hanno perso: una persona poteva capirne uno per intero. \prg è stato
progettato per restituire quella qualità, sull'architettura che gli studenti
incontreranno nell'industria (Tabella~\ref{tab:hist:c64}). Lo schermo ha la
stessa dimensione, e la stessa immediatezza, a quarant'anni di distanza.

\begin{table}[htb]
\centering\small
\caption{Un Commodore~64 e \prg.}\label{tab:hist:c64}
\begin{tabularx}{\textwidth}{@{}lLL@{}}
\toprule
 & \textbf{Commodore 64 (1982), MOS 6510} & \textbf{\prg su ESP32-C6, \riscv}\\
\midrule
Dimensione della parola & 8 bit & 32 bit\\
Registri & A, X, Y (più puntatore allo stack, stato, PC) & 32 registri di uso generale, da \reg{x0} a \reg{x31}\\
Clock & circa 1 MHz & 160 MHz\\
Spazio di indirizzamento & 64 KB & 4 GB\\
Insieme di istruzioni & 56 istruzioni, proprietario & una quarantina di istruzioni di base più le estensioni, standard aperto\\
Schermo & 320 per 200 pixel, 16 colori & viewport di 320 per 200, palette RGB565 di 256 voci\\
Si programma & scrivendo sulla macchina stessa & compilando una cartuccia e caricandola\\
\bottomrule
\end{tabularx}
\end{table}

Il 6502 è ancora un piacere da esplorare, per esempio con il simulatore su
\url{https://skilldrick.github.io/easy6502}.

\subsection{I giochi rétro come esemplari}

Un biologo studia il moscerino della frutta non perché i moscerini siano
importanti, ma perché un moscerino è abbastanza piccolo da essere capito per
intero e abbastanza ricco da mostrare come funziona l'ereditarietà. Un gioco
rétro svolge qui la stessa parte (Figura~\ref{fig:hist:games}): è abbastanza
piccolo da tenerlo in mente e abbastanza ricco da mettere in esercizio stato,
ingresso, aritmetica, decisioni e un ciclo, cioè l'intero corso. I quattro
giochi della figura sono cartucce di esempio del repository di \prg
(\fname{examples/games}).

\begin{figure}[htb]
\centering
\begin{subfigure}[b]{0.24\textwidth}\includegraphics[width=\textwidth]{game_pong.png}\caption{Pong}\end{subfigure}\hfill
\begin{subfigure}[b]{0.24\textwidth}\includegraphics[width=\textwidth]{game_breakout.png}\caption{Breakout}\end{subfigure}\hfill
\begin{subfigure}[b]{0.24\textwidth}\includegraphics[width=\textwidth]{game_tetris.png}\caption{Tetris}\end{subfigure}\hfill
\begin{subfigure}[b]{0.24\textwidth}\includegraphics[width=\textwidth]{game_asteroids.png}\caption{Asteroids}\end{subfigure}
\caption{I giochi rétro come esemplari.}
\label{fig:hist:games}
\end{figure}

% ----------------------------------------------------------------------------
\section{L'informatica come scienza applicata}

Dove si colloca l'informatica tra le scienze? Un esempio dal Dipartimento in
cui questo libro è stato scritto rende concreta la risposta
(Figura~\ref{fig:hist:chain}). Sul mare si forma una tempesta: un fenomeno
naturale. La fisica e la matematica, le scienze di base, lo descrivono con
delle equazioni. L'informatica, una scienza applicata, trasforma le equazioni
in algoritmi e simulazioni. L'ingegneria costruisce i supercalcolatori, le reti
e i sensori che le eseguono in modo affidabile. E l'economia riceve previsioni
e servizi che raggiungono le persone e creano valore. Le previsioni
meteorologiche e marine di \url{https://meteo.uniparthenope.it} sono calcolate
ogni giorno in questo modo, sui sistemi ad alte prestazioni del Dipartimento.

\begin{figure}[htb]
\centering
\begin{tikzpicture}[node distance=4mm]
  \node[blkacc,text width=22mm,minimum height=17mm] (a) {\textbf{Fenomeno naturale}\\[1pt]\footnotesize una tempesta};
  \node[blkamb,text width=22mm,minimum height=17mm,right=of a] (b) {\textbf{Scienze di base}\\[1pt]\footnotesize equazioni};
  \node[blkvio,text width=22mm,minimum height=17mm,right=of b] (c) {\textbf{Scienze applicate}\\[1pt]\footnotesize algoritmi};
  \node[blkgrn,text width=22mm,minimum height=17mm,right=of c] (d) {\textbf{Ingegneria}\\[1pt]\footnotesize sistemi};
  \node[blkred,text width=22mm,minimum height=17mm,right=of d] (e) {\textbf{Economia}\\[1pt]\footnotesize servizi};
  \foreach \x/\y in {a/b,b/c,c/d,d/e} \draw[->,thick] (\x) -- (\y);
\end{tikzpicture}
\caption{Dove si colloca l'informatica.}
\label{fig:hist:chain}
\end{figure}

Le stesse idee che in questo libro si esercitano su un microcontrollore a
160~MHz scalano fino ai cluster (Figura~\ref{fig:hist:hpc}): registri, memoria,
istruzioni, solo in numero molto maggiore.

\begin{figure}[htb]
\centering
\begin{subfigure}[b]{0.32\textwidth}\centering\includegraphics[width=\textwidth,height=36mm,keepaspectratio]{cluster_1.jpg}\caption{Un cluster didattico}\end{subfigure}\hfill
\begin{subfigure}[b]{0.32\textwidth}\centering\includegraphics[width=\textwidth,height=36mm,keepaspectratio]{cluster_2.jpg}\caption{Una simulazione parallela}\end{subfigure}\hfill
\begin{subfigure}[b]{0.32\textwidth}\centering\includegraphics[width=\textwidth,height=36mm,keepaspectratio]{rcf_page.png}\caption{\url{rcf.uniparthenope.it}}\end{subfigure}
\caption{Il calcolo ad alte prestazioni alla Parthenope: mini PC con GPU su una
rete a 10~Gbit, e le Research Computing Facilities.}
\label{fig:hist:hpc}
\end{figure}

% ----------------------------------------------------------------------------
\section{Verificate di aver capito}

\begin{quiz}{Quale idea è il cuore dell'architettura di von Neumann?}
\item Il processore ha esattamente quattro registri
\item Programma e dati sono memorizzati nella stessa memoria
\item Ogni istruzione richiede un ciclo di clock
\item Ingresso e uscita sono facoltativi
\end{quiz}

\begin{quiz}{Che cosa distingue \riscv da x86 e ARM?}
\item È l'insieme di istruzioni più veloce
\item È stata la prima architettura RISC
\item Il suo insieme di istruzioni è uno standard aperto, libero da realizzare
\item Funziona solo sui microcontrollori
\end{quiz}

\begin{quiz}{Chi è considerato il primo programmatore della storia?}
\item Charles Babbage
\item Alan Turing
\item Ada Lovelace
\item Konrad Zuse
\end{quiz}

\section{Esercizi}

\begin{exercise}[prima su carta, poi su \prg]
\begin{enumerate}[nosep]
\item Disegnate le quattro unità dell'architettura di von Neumann e scrivete,
  dentro ciascuna, una cosa del programma del Capitolo~\ref{ch:firstcontact}
  che vi risiede o vi accade.
\item Per ciascuna delle seguenti cose, dite se è un modo nuovo di
  \emph{rappresentare} l'informazione o un modo nuovo di \emph{prescrivere}
  operazioni: l'abaco, il telaio di Jacquard, la Pascalina, l'algebra di Boole,
  le schede perforate della macchina analitica, il programma memorizzato.
\item Scaricate dallo Store due dei quattro giochi della
  Figura~\ref{fig:hist:games}, oppure compilateli da \fname{examples/games}.
  Con i contatori del debugger, trovate quante istruzioni esegue ciascuno per
  frame, e confrontate i due giochi.
\end{enumerate}
\expected{Registri e ALU: lì è stata eseguita \ins{addi}. Memoria: il
contatore, e le istruzioni stesse. Unità di controllo: il prelievo di ogni
istruzione. Ingresso/uscita: la barra sullo schermo. Il programma memorizzato
appartiene a entrambe le classi: un programma è esso stesso informazione
rappresentata.}
\end{exercise}

\section{Glossario}

\begin{glossarytable}
computer science / informatics & informatica\\
computer & calcolatore, elaboratore\\
punched card & scheda perforata\\
loom & telaio\\
mainframe & mainframe, grande elaboratore\\
stored program & programma memorizzato\\
control unit & unit\`a di controllo\\
arithmetic logic unit & unit\`a aritmetico-logica\\
instruction set & insieme (set) di istruzioni\\
high-performance computing & calcolo ad alte prestazioni\\
\end{glossarytable}

\section{Conclusioni e letture}

La computazione è più antica dell'elettronica: rappresentazione e procedura
sono venute prima. Il calcolatore a programma memorizzato del 1945 è ancora la
macchina che si programma oggi. RISC ha reso i processori abbastanza semplici
da essere capiti per intero; \riscv li ha resi aperti. L'informatica è
pervasiva: l'informazione digitale appartiene a ogni campo, i dati ne sono la
moneta, e gli algoritmi trasformano i dati in decisioni.

\paragraph{Letture.} \citet{patterson2020}, capitolo~1; \emph{Playing with
RISC-V}, capitolo~1, «Why a Game, Why RISC-V, Why Now»
\citep{prg32textbook}; Charles Petzold, \emph{Code: The Hidden Language of
Computer Hardware and Software}; per la storia in profondità,
\citet{ceruzzi2003}.